Let $a=a_1<\cdots<a_s$ be positive integers with $a_{i+1}-a_i\le K$, where $a,K\ge1$. Must sufficiently large $a_s$, depending only on $a,K$, force equal sums over two distinct nonempty intervals of indices? Yes. We give an elementary proof with the explicit, deliberately nonoptimal choice
\[
 m=\left\lceil(a+K)^2e^{48K+10}\right\rceil,\qquad
 f(a,K)=a+K(2m-1).
 \tag{1}
\]
Hegyv\'ari proved the stronger bound $f(a,K)\ll ae^{O(K)}$. Our proof establishes the requested existence, but does not recover that sharper dependence on $a$.
